% ============================================================
% LAMPIRAN A - DAFTAR POKOK PERTANYAAN WAWANCARA
% Untuk dikirim ke Sekretariat ABW atas permintaan (19 Juni 2026)
% Pendamping proposal-anies.tex. Kompilasi: pdflatex (Overleaf-ready).
% Disusun ulang 20 Juni 2026 (revisi): 7 pertanyaan terbuka dan generatif (metode funnel);
% pendalaman dilakukan lewat probe netral saat wawancara. Netral, tidak mengarahkan, bukan kritik.
% Lihat DECISIONS.md.
% ============================================================
\documentclass[11pt,a4paper]{article}

\usepackage[T1]{fontenc}
\usepackage[utf8]{inputenc}
\usepackage[indonesian]{babel}
\usepackage[margin=2.5cm]{geometry}
\usepackage{setspace}
\usepackage{enumitem}
\usepackage{parskip}
\usepackage[hidelinks]{hyperref}

\setstretch{1.15}

\begin{document}

% ---------- KOP ----------
\begin{center}
  {\Large\bfseries LAMPIRAN A}\\[4pt]
  {\large Daftar Pokok Pertanyaan Wawancara}\\[2pt]
  {\normalsize Belajar tentang Pengalaman Berjalan Kaki di Jakarta}\\[2pt]
  {\small Narasumber: Bapak H.~Anies Rasyid Baswedan, Ph.D.}\\
  {\small Tim beranggotakan lima peserta, Apple Developer Academy @ BINUS}
\end{center}

\vspace{4pt}
\hrule
\vspace{12pt}

\noindent
Wawancara bersifat \textbf{semi-terstruktur} dan diperkirakan berlangsung sekitar 60 menit.
Pertanyaan di bawah ini bersifat \textbf{terbuka} dan menjadi pokok bahasan; kami berharap dapat
belajar dari pandangan, pengalaman, dan gagasan Bapak. Percakapan dapat berkembang secara alami,
dan kami akan menggali lebih dalam mengikuti arah cerita Bapak. Sebagian tema kami kaitkan dengan
ketertarikan kami pada pengalaman berjalan kaki di Jakarta.

\vspace{6pt}

% ---------- BAGIAN 1 ----------
\section*{Bagian 1: Filosofi dan Pengalaman}
\begin{enumerate}[label=\arabic*., itemsep=6pt, leftmargin=*]
  \item Bagaimana Bapak memandang peran berjalan kaki dalam kehidupan sebuah kota?
  \item Menurut Bapak, apa yang membuat orang merasa nyaman dan betah berjalan kaki di sebuah
        kawasan?
  \item Bisa Bapak ceritakan pengalaman yang paling berkesan saat menata ruang pejalan kaki di
        Jakarta?
\end{enumerate}

% ---------- BAGIAN 2 ----------
\section*{Bagian 2: Warga dan Kota}
\begin{enumerate}[label=\arabic*., itemsep=6pt, leftmargin=*, start=4]
  \item Bagaimana Bapak melihat peran warga dalam ikut membangun dan merawat kotanya?
  \item Dari pengalaman Bapak, apa yang membuat suara dan partisipasi warga dapat benar-benar
        membawa perubahan bagi kota?
\end{enumerate}

% ---------- BAGIAN 3 ----------
\section*{Bagian 3: Menjangkau Semua dan Masa Depan}
\begin{enumerate}[label=\arabic*., itemsep=6pt, leftmargin=*, start=6]
  \item Bagaimana Bapak membayangkan agar kualitas ruang pejalan kaki yang baik dapat tumbuh dan
        menjangkau makin banyak lingkungan tempat warga tinggal?
  \item Bagaimana Bapak membayangkan masa depan kota yang ramah dan manusiawi bagi pejalan kaki,
        dan apa pesan Bapak bagi kami yang ingin ikut mewujudkannya?
\end{enumerate}

\vspace{10pt}
\hrule
\vspace{8pt}
\noindent\small\textit{Catatan etika:} hasil wawancara digunakan hanya untuk pembelajaran dan
pengembangan solusi non-komersial di lingkup tim kami. Perekaman audio hanya atas izin Bapak, dan
transkrip maupun kutipan dapat kami kirimkan untuk Bapak tinjau sebelum digunakan.

\end{document}
